\bookchapter{The Shell: Talking to Your Computer in Words}{ch:02-the-shell-talking-to-your-computer-in-words}

\begin{chapterintro}

This chapter turns the shell from a mystery into a tool. You'll learn to read any command, move around, see how the shell finds programs, and join small commands into bigger ones. By the end, you'll predict what a command does before you press Enter.

\end{chapterintro}

\emph{The problem.} I type a command and something happens, but I couldn't tell you what.

\emph{The question.} What exactly happens between pressing Enter and seeing the result?

\section{The Shell's Loop}\label{02-the-shell-talking-to-your-computer-in-words:the-shells-loop}

The shell does the same four things, forever:

\begin{enumerate}
\def\labelenumi{\arabic{enumi}.}
\tightlist
\item
  It shows a prompt, the text before your cursor (often your name, the folder you're in, and a \texttt{\$} or \texttt{\%}).
\item
  It reads the line you type when you press Enter.
\item
  It finds the program you named and runs it, passing along the rest of the line.
\item
  It shows the program's output, then goes back to step 1.
\end{enumerate}

The two shells you'll meet are bash, the default on most Linux systems, and zsh, the default on macOS. For everything in this book they behave the same way.

\section{The Parts of a Command}\label{02-the-shell-talking-to-your-computer-in-words:the-parts-of-a-command}

Every command line has the same shape: a command, then options, then arguments.

\codeneed{5}\begin{codeblock}
\begin{Verbatim}[breaklines=false, fontsize=\fontsize{8.8}{8.68}\selectfont, fontfamily=diagrammono, formatcom=\diagramlines]
  ls   -l   -a   /home/ada/sandbox
  │    └──┬──┘   └──────┬────────┘
  │       │             └─ argument: what to act on
  │       └─ options: how to do it
  └─ command: which program to run
\end{Verbatim}
\end{codeblock}

\begin{itemize}
\tightlist
\item
  The command is the program to run. Here it's \texttt{ls}, which lists files.
\item
  An \textbf{option}\index{option} (also called a flag) changes how the command behaves. Short options are one letter after a dash (\texttt{-l}), and can be combined (\texttt{-la} means \texttt{-l\ -a}). Long options are words after two dashes (\texttt{-\/-all}).
\item
  An \textbf{argument}\index{argument} is the thing the command acts on, usually a file or folder. Most commands have a sensible default when you leave it out.
\end{itemize}

Spaces separate the parts. That's why file names with spaces cause trouble: \texttt{my\ notes.txt} looks like two arguments. Wrap such names in quotes: \texttt{"my\ notes.txt"}.

\section{Where You Are: Paths and the Working Directory}\label{02-the-shell-talking-to-your-computer-in-words:where-you-are-paths-and-the-working-directory}

The shell always has a current folder, called the \textbf{working directory}\index{working directory}. Commands act there unless you say otherwise. When you open a terminal, you start in your \textbf{home directory}\index{home directory}, your personal folder: \texttt{/home/ada} on Linux, \texttt{/Users/ada} on macOS. The shell lets you write it as \texttt{\textasciitilde{}}.

A \textbf{path}\index{path} is the address of a file or folder. An absolute path starts at the very top of the file system, \texttt{/}, and works from anywhere: \texttt{/}\allowbreak{}\texttt{home/}\allowbreak{}\texttt{ada/}\allowbreak{}\texttt{sandbox/}\allowbreak{}\texttt{todo.}\allowbreak{}\texttt{txt}. A relative path starts from where you are now: if you're in \texttt{/home/ada}, then \texttt{sandbox/todo.txt} means the same file. Two special names help: \texttt{.} means ``this folder'' and \texttt{..} means ``the folder above''.

Three commands do most of the moving around. \texttt{pwd} prints the working directory. \texttt{ls} lists what's in it. \texttt{cd} changes it.

\codeneed{8}\begin{codeblock}

\begin{Shaded}
\begin{Highlighting}[]
\BuiltInTok{pwd}
\FunctionTok{mkdir} \AttributeTok{{-}p}\NormalTok{ sandbox/notes}
\BuiltInTok{cd}\NormalTok{ sandbox}
\FunctionTok{touch}\NormalTok{ todo.txt ideas.txt}
\FunctionTok{ls}
\end{Highlighting}
\end{Shaded}

\end{codeblock}

\codeneed{2}\begin{codeblock}
\begin{Verbatim}[formatcom=\outputstyle]
/home/ada
ideas.txt  notes  todo.txt
\end{Verbatim}
\end{codeblock}

\codeneed{9}

\texttt{mkdir\ -p} made the folder \texttt{sandbox} and, inside it, \texttt{notes} (the \texttt{-p} means ``make parent folders too, and don't complain if they exist''). \texttt{touch} created two empty files. Now the long listing:

\codeneed{6}\begin{codeblock}

\begin{Shaded}
\begin{Highlighting}[]
\FunctionTok{ls} \AttributeTok{{-}l}
\end{Highlighting}
\end{Shaded}

\end{codeblock}

\codeneed{4}\begin{codeblock}
\begin{Verbatim}[formatcom=\outputstyle]
total 4
-rw-r--r-- 1 ada ada    0 Sep 26 14:27 ideas.txt
drwxr-xr-x 2 ada ada 4096 Sep 26 14:27 notes
-rw-r--r-- 1 ada ada    0 Sep 26 14:27 todo.txt
\end{Verbatim}
\end{codeblock}

Each line shows permissions, owner, size in bytes, date, and name. A \texttt{d} at the start marks a directory. Chapter 3 reads the rest.

To move around, \texttt{cd\ notes} goes into a folder, \texttt{cd\ ..} climbs one level, and \texttt{cd} on its own always takes you home.

\Needspace{16\baselineskip}

\section{Making, Copying, and Removing}\label{02-the-shell-talking-to-your-computer-in-words:making-copying-and-removing}

A few more commands handle everyday file work:

\begin{booktable}
\begin{longtable}{>{\raggedright\arraybackslash}p{\dimexpr 0.2623\linewidth-1.4754\tabcolsep\relax}>{\raggedright\arraybackslash}p{\dimexpr 0.7377\linewidth-0.5246\tabcolsep\relax}}
\tabletop
\rowcolor{tablehead}\thead{Command} & \thead{What it does} \\
\tablemid
\endfirsthead
\tabletop
\rowcolor{tablehead}\thead{Command} & \thead{What it does} \\
\tablemid
\endhead
\tablebottom
\endlastfoot
\texttt{mkdir\ -p\ a/b} & Make folders, including parents \\
\texttt{touch\ f} & Create an empty file (or update its time) \\
\texttt{cp\ src\ dst} & Copy a file; \texttt{cp\ -r} copies a folder \\
\texttt{mv\ old\ new} & Move or rename \\
\texttt{rm\ f} & Remove a file; \texttt{rm\ -r} removes a folder \\
\texttt{cat\ f} & Print a whole file \\
\end{longtable}
\end{booktable}

\begin{callout}{warning}

\texttt{rm} does not use the bin or the Trash. The file is gone the moment you press Enter. Be especially careful with \texttt{rm\ -r}, and read the whole line before you press Enter.

\end{callout}

\section{How the Shell Finds a Command}\label{02-the-shell-talking-to-your-computer-in-words:how-the-shell-finds-a-command}

When you type \texttt{ls}, how does the shell know which program you mean? It looks in a list of folders stored in a setting called \textbf{PATH}\index{path@PATH}. It checks them from left to right and runs the first program with that name.

\codeneed{3}\begin{codeblock}

\begin{Shaded}
\begin{Highlighting}[]
\BuiltInTok{echo} \VariableTok{$PATH}
\end{Highlighting}
\end{Shaded}

\end{codeblock}

\codeneed{1}\begin{codeblock}
\begin{Verbatim}[formatcom=\outputstyle]
/usr/local/bin:/usr/bin:/bin:/usr/local/games:/usr/games
\end{Verbatim}
\end{codeblock}

\codeneed{7}

The folders are separated by colons. \texttt{which} tells you where a command was found:

\codeneed{4}\begin{codeblock}

\begin{Shaded}
\begin{Highlighting}[]
\FunctionTok{which}\NormalTok{ ls python3}
\end{Highlighting}
\end{Shaded}

\end{codeblock}

\codeneed{2}\begin{codeblock}
\begin{Verbatim}[formatcom=\outputstyle]
/usr/bin/ls
/usr/local/bin/python3
\end{Verbatim}
\end{codeblock}

\codeneed{6}

A few commands, like \texttt{cd}, aren't programs at all. They're a builtin, part of the shell itself:

\codeneed{3}\begin{codeblock}

\begin{Shaded}
\begin{Highlighting}[]
\BuiltInTok{type}\NormalTok{ cd}
\end{Highlighting}
\end{Shaded}

\end{codeblock}

\codeneed{1}\begin{codeblock}
\begin{Verbatim}[formatcom=\outputstyle]
cd is a shell builtin
\end{Verbatim}
\end{codeblock}

\codeneed{4}

This explains the most common error in the terminal:

\codeneed{1}\begin{codeblock}
\begin{Verbatim}[formatcom=\outputstyle]
bash: nosuchcmd: command not found
\end{Verbatim}
\end{codeblock}

No folder on your PATH holds a program with that name. It's misspelled, not installed, or installed in a folder that isn't on your PATH.

\section{Streams, Redirection, and Pipes}\label{02-the-shell-talking-to-your-computer-in-words:streams-redirection-and-pipes}

Every program starts with three streams, like three hoses attached to it. \textbf{Standard input}\index{standard input} (stdin) is where it reads from, usually your keyboard. \textbf{Standard output}\index{standard output} (stdout) is where it writes results, usually your screen. \textbf{Standard error}\index{standard error} (stderr) is a second output just for error messages, also your screen.

\textbf{Redirection}\index{redirection} points those hoses somewhere else. \texttt{\textgreater{}} sends stdout into a file, replacing what was there. \texttt{\textgreater{}\textgreater{}} adds to the end instead. \texttt{2\textgreater{}} sends stderr into a file.

\codeneed{7}\begin{codeblock}

\begin{Shaded}
\begin{Highlighting}[]
\BuiltInTok{cd}\NormalTok{ \textasciitilde{}/sandbox}
\BuiltInTok{echo} \StringTok{"buy milk"} \OperatorTok{\textgreater{}}\NormalTok{ todo.txt}
\BuiltInTok{echo} \StringTok{"call Sam"} \OperatorTok{\textgreater{}\textgreater{}}\NormalTok{ todo.txt}
\FunctionTok{cat}\NormalTok{ todo.txt}
\end{Highlighting}
\end{Shaded}

\end{codeblock}

\codeneed{2}\begin{codeblock}
\begin{Verbatim}[formatcom=\outputstyle]
buy milk
call Sam
\end{Verbatim}
\end{codeblock}

\codeneed{9}

Errors travel on their own hose, so you can separate them:

\codeneed{6}\begin{codeblock}

\begin{Shaded}
\begin{Highlighting}[]
\FunctionTok{ls}\NormalTok{ /nope}
\FunctionTok{ls}\NormalTok{ /nope }\DecValTok{2}\OperatorTok{\textgreater{}}\NormalTok{ errors.txt}
\FunctionTok{cat}\NormalTok{ errors.txt}
\end{Highlighting}
\end{Shaded}

\end{codeblock}

\codeneed{2}\begin{codeblock}
\begin{Verbatim}[formatcom=\outputstyle]
ls: cannot access '/nope': No such file or directory
ls: cannot access '/nope': No such file or directory
\end{Verbatim}
\end{codeblock}

The first \texttt{ls} printed its complaint on the screen. The second printed nothing, because its error went into \texttt{errors.txt}. On macOS the message reads \texttt{ls:}\allowbreak{}\texttt{\ /}\allowbreak{}\texttt{nope:}\allowbreak{}\texttt{\ No\ such\ file\ or\ directory}; the idea is the same.

\codeneed{6}

A \textbf{pipe}\index{pipe}, written \texttt{\textbar{}}, connects one program's stdout to the next program's stdin. This is how small tools join forces. How many entries are in \texttt{/etc}? \texttt{ls} lists them, \texttt{wc\ -l} counts lines:

\codeneed{3}\begin{codeblock}

\begin{Shaded}
\begin{Highlighting}[]
\FunctionTok{ls}\NormalTok{ /etc }\KeywordTok{|} \FunctionTok{wc} \AttributeTok{{-}l}
\end{Highlighting}
\end{Shaded}

\end{codeblock}

\codeneed{1}\begin{codeblock}
\begin{Verbatim}[formatcom=\outputstyle]
75
\end{Verbatim}
\end{codeblock}

Read a pipeline left to right, like a sentence: ``list \texttt{/etc}, then count the lines''. Chapter 5 builds longer ones.

\section{Getting Help}\label{02-the-shell-talking-to-your-computer-in-words:getting-help}

You never need to memorise options. Every standard command has a manual page, or \textbf{man page}\index{man page}. Run \texttt{man\ ls} to read it; use the arrow keys or space to scroll, \texttt{/word} to search, and \texttt{q} to quit. Most GNU tools on Linux also print a short summary with \texttt{-\/-help}:

\codeneed{3}\begin{codeblock}

\begin{Shaded}
\begin{Highlighting}[]
\FunctionTok{ls} \AttributeTok{{-}{-}help} \KeywordTok{|} \FunctionTok{head} \AttributeTok{{-}1}
\end{Highlighting}
\end{Shaded}

\end{codeblock}

\codeneed{1}\begin{codeblock}
\begin{Verbatim}[formatcom=\outputstyle]
Usage: ls [OPTION]... [FILE]...
\end{Verbatim}
\end{codeblock}

macOS's \texttt{ls} doesn't understand \texttt{-\/-help}; use \texttt{man\ ls} there.

\begin{callout}{tip}

Press Tab while typing a name, and the shell completes it for you; press it twice to see the choices. It saves typing and, more importantly, typos.

\end{callout}

A few more keys make life easier. The Up arrow brings back earlier commands. Ctrl-R searches your history as you type. Ctrl-C cancels whatever is running. Ctrl-L clears the screen.

\section{Lab: A Sandbox Tour}\label{02-the-shell-talking-to-your-computer-in-words:lab-a-sandbox-tour}

Work in a new terminal window. Each step says what you should see.

\begin{enumerate}
\def\labelenumi{\arabic{enumi}.}
\tightlist
\item
  Make a fresh folder for this lab: \texttt{mkdir\ -}\allowbreak{}\texttt{p\ \textasciitilde{}/}\allowbreak{}\texttt{sandbox/}\allowbreak{}\texttt{ch2/}\allowbreak{}\texttt{notes}, then \texttt{cd\ \textasciitilde{}/sandbox/ch2}. \textbf{Expected:} \texttt{pwd} prints your home folder followed by \texttt{/sandbox/ch2}.
\item
  Create a file with two lines: \texttt{echo\ "first"\ \textgreater{}\ log.txt}, then \texttt{echo\ "second"\ \textgreater{}\textgreater{}\ log.}\allowbreak{}\texttt{txt}, then \texttt{cat\ log.txt}. \textbf{Expected:} \texttt{first} and \texttt{second} on two lines.
\item
  Replace instead of adding: \texttt{echo\ "only"\ \textgreater{}\ log.txt} and \texttt{cat\ log.txt}. \textbf{Expected:} just \texttt{only}. That's the difference between \texttt{\textgreater{}} and \texttt{\textgreater{}\textgreater{}}.
\item
  Copy and move: \texttt{cp\ log.}\allowbreak{}\texttt{txt\ notes/}\allowbreak{}\texttt{copy.}\allowbreak{}\texttt{txt}, then \texttt{mv\ notes/}\allowbreak{}\texttt{copy.}\allowbreak{}\texttt{txt\ notes/}\allowbreak{}\texttt{moved.}\allowbreak{}\texttt{txt}, then \texttt{ls\ notes}. \textbf{Expected:} \texttt{moved.txt}.
\item
  Separate an error: \texttt{ls\ /}\allowbreak{}\texttt{nope\ 2\textgreater{}\ err.txt}, then \texttt{cat\ err.txt}. \textbf{Expected:} nothing on screen from the first command; the error message from the second.
\item
  Count with a pipe: \texttt{ls\ \textasciitilde{}/}\allowbreak{}\texttt{sandbox/}\allowbreak{}\texttt{ch2\ \textbar{}\ wc\ -}\allowbreak{}\texttt{l}. \textbf{Expected:} \texttt{3} (\texttt{err.txt}, \texttt{log.txt}, \texttt{notes}).
\item
  Find where Python lives: \texttt{which\ python3}. \textbf{Expected:} a path such as \texttt{/usr/bin/python3}.
\item
  Look up an option: run \texttt{man\ wc} and find what \texttt{-w} does. \textbf{Expected:} it counts words. Press \texttt{q} to leave.
\end{enumerate}

\section{Summary, Key Terms, and Review Questions}\label{02-the-shell-talking-to-your-computer-in-words:summary-key-terms-and-review-questions}

\subsection*{Summary}\label{02-the-shell-talking-to-your-computer-in-words:summary}

\begin{itemize}
\tightlist
\item
  The shell loops: prompt, read, run, show. bash and zsh work the same way for everyday use.
\item
  A command line is a command, then options, then arguments, separated by spaces.
\item
  The working directory is where commands act. Absolute paths start at \texttt{/}; relative paths start where you are.
\item
  The shell finds programs by searching the folders on your PATH, left to right.
\item
  Every program has stdin, stdout, and stderr. Redirection points them at files; pipes connect programs.
\item
  \texttt{man} and \texttt{-\/-help} explain any command.
\end{itemize}

\Needspace{17\baselineskip}

\subsection*{Key Terms}\label{02-the-shell-talking-to-your-computer-in-words:key-terms}

\begin{booktable}
\begin{longtable}{>{\raggedright\arraybackslash}p{\dimexpr 0.3667\linewidth-1.2667\tabcolsep\relax}>{\raggedright\arraybackslash}p{\dimexpr 0.6333\linewidth-0.7333\tabcolsep\relax}}
\tabletop
\rowcolor{tablehead}\thead{Term} & \thead{Meaning} \\
\tablemid
\endfirsthead
\tabletop
\rowcolor{tablehead}\thead{Term} & \thead{Meaning} \\
\tablemid
\endhead
\tablebottom
\endlastfoot
Option & A setting that changes how a command behaves, like \texttt{-l} \\
Argument & What a command acts on, usually a file or folder \\
Working directory & The folder the shell is currently in \\
Home directory & Your personal folder, written \texttt{\textasciitilde{}} \\
Path & The address of a file or folder \\
PATH & The list of folders the shell searches for commands \\
Standard input, output, error & The three streams every program starts with \\
Redirection & Sending a stream to or from a file (\texttt{\textgreater{}},
\texttt{\textgreater{}\textgreater{}}, \texttt{2\textgreater{}}) \\
Pipe & The ` \\
Man page & The built-in manual for a command \\
\end{longtable}
\end{booktable}

\subsection*{Review Questions}\label{02-the-shell-talking-to-your-computer-in-words:review-questions}

\begin{enumerate}
\def\labelenumi{\arabic{enumi}.}
\tightlist
\item
  In \texttt{cp\ -}\allowbreak{}\texttt{r\ photos\ /}\allowbreak{}\texttt{home/}\allowbreak{}\texttt{ada/}\allowbreak{}\texttt{backup}, which part is the command, which the option, and which the arguments?
\item
  You're in \texttt{/home/ada/sandbox}. Write the absolute and relative paths to \texttt{notes/ideas.txt}.
\item
  What's the difference between \texttt{\textgreater{}} and \texttt{\textgreater{}\textgreater{}}? Which one can destroy data?
\item
  What does ``command not found'' tell you, and what are three possible causes?
\item
  What does \texttt{ls\ /etc\ \textbar{}\ wc\ -l} do, step by step?
\end{enumerate}

\begin{understandbox}

\textbf{You understand this when} you can read a command like \texttt{ls\ -}\allowbreak{}\texttt{l\ \textasciitilde{}/}\allowbreak{}\texttt{sandbox\ 2\textgreater{}\ err.}\allowbreak{}\texttt{txt\ \textbar{}\ wc\ -}\allowbreak{}\texttt{l} aloud and say what each part does.

\end{understandbox}
